\documentclass[11pt]{article}
\usepackage[margin=1in]{geometry}
\usepackage{amsmath, amssymb, amsthm}
\usepackage{booktabs}
\usepackage{enumitem}

\newcommand{\mesh}{\mathrm{mesh}}
\newcommand{\SS}{\mathrm{SS}}

\title{Linearization in Microdata Generation}
\date{}

\begin{document}
\maketitle

\section{Overview}

The microdata simulation produces household-level distributions at each period~$t$
by combining three objects:
\begin{enumerate}[label=(\roman*)]
    \item \textbf{Distribution} $\mu_t(b,k,h)$: joint probability mass over the grid,
    \item \textbf{Policies} $x^a_t(b,k,h),\; x^n_t(b,k,h)$: savings/consumption decisions,
    \item \textbf{Income} $y_t(b,k,h)$: disposable income at each grid point.
\end{enumerate}
Each of these involves a separate linearization decision.
The choices are independent and summarized in Table~\ref{tab:summary}.

\begin{table}[h]
\centering
\begin{tabular}{@{}llll@{}}
\toprule
Object & Treatment & Reason & Section \\
\midrule
Distribution $\mu_t$ & Smoothed perturbation & Softplus positivity & \ref{sec:distribution} \\
Policies $x^a_t, x^n_t$ & Linearized (JVP) & Forward-looking, requires EGM & \ref{sec:policies} \\
Nonasset income & Nonlinear & Static mapping, exact & \ref{sec:nonasset} \\
Rental income & Dynamic $(R^K_t)$ & Aggregate scaling, well-behaved & \ref{sec:asset} \\
Liquid return & Fixed at SS & Discontinuity at $b=0$ & \ref{sec:asset} \\
\bottomrule
\end{tabular}
\caption{Linearization choices in the microdata simulation.}
\label{tab:summary}
\end{table}

\section{Aggregate State Construction}

All aggregate variables entering the household problem are stored in a vector
$\mathbf{a}_t = (w_t H_t,\, N_t,\, R^K_t,\, R^L_t,\, R^D_t,\, \ldots)$.
The linear state-space representation produces log-deviations from steady state:
\begin{equation}\label{eq:lom}
    s_{t+1} = \Phi\, s_t + B\, \varepsilon_t, \qquad
    c_t = \Theta\, s_t,
\end{equation}
where $s_t$ are states (including exogenous shocks), $c_t$ are controls,
$\Phi$ is the law of motion, and $\Theta$ maps states to controls.
The aggregate variables are recovered as \emph{levels}:
\begin{equation}\label{eq:exp}
    a_{t,j} = \exp\!\bigl(\log a^{\SS}_j + \delta_{t,j}\bigr)
            = a^{\SS}_j \cdot e^{\delta_{t,j}},
\end{equation}
where $\delta_{t,j}$ is the log-deviation from \eqref{eq:lom}.
This exponential transformation is the source of nonlinearity in
an otherwise linear simulation.

\section{Distribution}\label{sec:distribution}

The joint distribution is reconstructed from perturbed marginals and a copula.
Marginal PDFs are perturbed around their steady-state values using
a softplus-based positivity mapping:
\begin{equation}
    \tilde{p}_i = \frac{\operatorname{softplus}(p^{\SS}_i + \delta_i;\, \kappa) + p_{\min}}
                       {\sum_j \bigl[\operatorname{softplus}(p^{\SS}_j + \delta_j;\, \kappa) + p_{\min}\bigr]},
\end{equation}
where $\operatorname{softplus}(x;\kappa) = \kappa^{-1}\log(1+e^{\kappa x})$ is a smooth
approximation to $\max(x,0)$, $\kappa$ controls sharpness, and $p_{\min}$
provides a floor. This replaces the hard clamp $\max(\cdot, 0)$ which created
discrete on/off switching of borrowing-constraint mass and moment jumps.

\section{Policies}\label{sec:policies}

Consumption and savings policies $x^a_t(b,k,h)$ and $x^n_t(b,k,h)$ are
\emph{forward-looking}: they depend on expected future marginal values through
the Euler equation and are solved via the Endogenous Grid Method (EGM).
The EGM maps
\[
    (x^a_t, x^n_t)
    = \mathcal{E}\!\bigl(\mathbb{E}_t[W'_{b,t+1}],\;
                         \mathbb{E}_t[W'_{k,t+1}],\;
                         \mathbf{a}_t,\;
                         \mathbf{y}_t\bigr),
\]
where $W'_{b}, W'_{k}$ are marginal continuation values and $\mathcal{E}$
denotes the full EGM operator (nonlinear, involving interpolation, constraints,
and max operators).

Because $\mathcal{E}$ is expensive and nonlinear, we linearize via a
\emph{Jacobian--vector product} (JVP) using forward-mode automatic differentiation
(ForwardDiff.jl). Writing all EGM inputs collectively as
$\mathbf{z}_t = (W'_{b,t},\, W'_{k,t},\, \mathbf{a}_t)$, the first-order
approximation is
\begin{equation}
    x_t \approx x^{\SS} + \nabla_{\mathbf{z}} \mathcal{E}\big|_{\SS}
                           \cdot (\mathbf{z}_t - \mathbf{z}^{\SS}).
\end{equation}
The JVP computes $\nabla_{\mathbf{z}} \mathcal{E} \cdot \Delta\mathbf{z}$
in a single forward pass by seeding
$\mathbf{z}^{\text{dual}} = \mathbf{z}^{\SS} + \epsilon\,\Delta\mathbf{z}_t$
and extracting the $\epsilon$-coefficient from the output.

\section{Nonasset Income}\label{sec:nonasset}

Nonasset income---net labor income, union profits, transfers, and entrepreneurial
profits---is a \emph{static} function of current-period aggregates and the
household's grid position:
\begin{equation}
    y^{\text{nonasset}}_t(b,k,h) = f\!\bigl(\mathbf{a}_t,\; h\bigr).
\end{equation}
It does not depend on expectations or continuation values. Therefore we evaluate
$f$ at the \emph{nonlinear} (level) aggregates $\mathbf{a}_t$ from \eqref{eq:exp}.
This is exact---no linearization is needed or desired.

An earlier implementation linearized income via ForwardDiff alongside the policies,
but this was abandoned: the income function contains discontinuities
(e.g., the borrower/saver indicator $\mathbf{1}_{b > 0}$ in the effective interest
rate $R^i$) that produce domain issues (NaN) in the automatic differentiation.
Since income is a known closed-form mapping, evaluating it directly is both
simpler and more accurate.

\section{Asset Income}\label{sec:asset}

Asset income combines rental income from illiquid assets and net interest on
liquid assets:
\begin{equation}\label{eq:asset_nl}
    y^{\text{asset}}_t(b,k,h) = \underbrace{(R^K_t - 1)\, k}_{\text{rental}}
                               \;+\; \underbrace{(R^i_t - 1)\, b}_{\text{liquid interest}},
\end{equation}
where $R^i_t = R^L_t \cdot \mathbf{1}_{b>0} + R^D_t \cdot \mathbf{1}_{b\leq 0}$.

The two components behave very differently:

\paragraph{Rental income: dynamic (no issue).}
Rental income $(R^K_t - 1) \cdot k$ responds to aggregate shocks, but the
response is \emph{purely aggregate}: every household's rental income scales by
the same factor $(R^K_t - 1) / (R^{K,\SS} - 1)$.  Relative to the
cross-sectional mean, the quintile ratios are therefore constant---rental income
generates no distributional dynamics beyond a level shift.  Because this
component is well-behaved (no discontinuities, no regime switching), we keep it
dynamic:
\begin{equation}
    y^{\text{rental}}_t(b,k,h) = (R^K_t - 1)\, k.
\end{equation}

\paragraph{Liquid return: fixed at steady state.}
The effective rate $R^i_t$ contains a structural discontinuity at $b = 0$:
savers receive $R^L_t$ while borrowers pay $R^D_t$.  In the baseline
calibration, $R^{L,\SS} = R^{RB} = 1.0$, so the net liquid return for savers
is \emph{exactly zero} at steady state.  Any perturbation to $R^L_t$---however
small---switches savers from zero to nonzero liquid return, creating a discrete
regime change that propagates into quintile means as binary switching patterns.
We therefore fix liquid return at its steady-state value:
\begin{equation}
    y^{\text{liquid}}_t(b,k,h) \equiv (R^{i,\SS} - 1)\, b.
\end{equation}

\paragraph{Combined asset income.}
\begin{equation}\label{eq:asset_fix}
    y^{\text{asset}}_t(b,k,h)
        = \underbrace{(R^K_t - 1)\, k}_{\text{dynamic}}
        \;+\; \underbrace{(R^{i,\SS} - 1)\, b}_{\text{fixed at SS}}.
\end{equation}

Note that the \texttt{incomes!()} function is not modified. Asset income is
recomputed from rates after the call to \texttt{incomes!()}, and the
rental and liquid components are stored separately for diagnostic purposes.

\section{Summary}

The term ``linearized'' in the microdata simulation refers to three distinct
choices:
\begin{enumerate}
    \item \textbf{Aggregate dynamics} (always): the state-space \eqref{eq:lom}
          is a first-order perturbation of the full nonlinear model.
    \item \textbf{Policies} (when \texttt{policy\_mode = :linearized}):
          the EGM operator is linearized via JVP around steady state.
    \item \textbf{Asset income}: rental income $(R^K_t - 1)\cdot k$ is kept
          dynamic; liquid return $(R^i - 1)\cdot b$ is fixed at steady state
          due to a structural discontinuity at $b=0$.
\end{enumerate}
Nonasset income and the distribution reconstruction are \emph{not} linearized;
the former is evaluated exactly, and the latter uses a smooth nonlinear
positivity mapping (softplus).

\end{document}
